\documentclass[preprint,12pt]{elsarticle}
\usepackage[english]{babel}
\usepackage[utf8]{inputenc}
\usepackage[T1]{fontenc}
\usepackage{amsmath}
\usepackage{amsthm}
\usepackage{amssymb}
\usepackage{graphicx}
\usepackage{bbm}
\usepackage{mathtools}
\usepackage{booktabs}
\usepackage{multirow}
\usepackage{hyperref}

% Notation macros (consistent with parent paper)
\newcommand{\RV}[1]{\mathrm{#1}}
\newcommand{\Prob}[1]{\mathrm{P}\{#1\}}
\newcommand{\Expect}[1]{\mathrm{E}\{#1\}}
\newcommand{\cv}{\RV{T}}           % Component lifetime
\newcommand{\sv}{\RV{S}}           % System lifetime
\newcommand{\cand}{\RV{C}}         % Candidate set
\newcommand{\tparam}[1]{{#1}^{\star}}  % True parameter
\newcommand{\estimator}[1]{\hat{#1}}   % MLE estimator
\newcommand{\infomatrx}{\mathcal{I}}
\newcommand{\pdf}{f}
\newcommand{\expdist}{\mathrm{EXP}}
\renewcommand{\vec}[1]{\boldsymbol{#1}}

% Theorem environments
\newtheorem{theorem}{Theorem}[section]
\newtheorem{corollary}[theorem]{Corollary}
\newtheorem{proposition}[theorem]{Proposition}
\theoremstyle{definition}
\newtheorem{definition}[theorem]{Definition}
\theoremstyle{remark}
\newtheorem{remark}[theorem]{Remark}

\hypersetup{
  colorlinks=true,
  linkcolor=black,
  citecolor=black,
  urlcolor=blue
}

\begin{document}

\begin{frontmatter}

\title{Information Recovery under Deterministic Masking\\in Exponential Series Systems}

\author[siue]{Alexander Towell}
\ead{lex@metafunctor.com}
\ead[url]{https://orcid.org/0000-0001-6443-9897}

\affiliation[siue]{organization={Southern Illinois University Edwardsville}}

\begin{abstract}
In masked system failure data, the candidate set of possibly failed components is
conventionally assumed to degrade estimation precision as its cardinality grows.
We show that for exponential series systems, a deterministic masking mechanism
that injectively maps each failure cause to a unique candidate set recovers the
complete-identification Fisher information matrix---regardless of the candidate
set cardinality. Information loss from masking is thus a property of the masking
\emph{mechanism}, not the masking \emph{cardinality}. We verify this result
numerically via Monte Carlo simulation.
\end{abstract}

\begin{keyword}
series systems \sep masked failure data \sep Fisher information \sep
deterministic masking \sep exponential distribution \sep competing risks
\MSC[2020] 62N05 \sep 62F12 \sep 62B10
\end{keyword}

\end{frontmatter}

\section{Introduction}
\label{sec:introduction}

In series (or $m$-out-of-$m$) reliability systems, all components must function
for the system to operate. When a system fails, diagnostics may identify only a
\emph{candidate set} $C \subseteq \{1,\ldots,m\}$ of components that could have
caused the failure, rather than the exact failure cause. This is known as
\emph{masked failure data}~\cite{miyakawa1984}.

A central question in the masked-data literature is how the masking mechanism
affects estimation precision. The conventional wisdom, supported by results on
mutual information between the failure cause and the candidate set, is that
larger candidate sets carry less information and therefore yield less precise
parameter estimates. Under the widely studied uniform masking
model~\cite{miyakawa1984, towell2023reliability, towell2025masked}, this monotonicity indeed
holds: the Fisher information matrix (FIM) degrades strictly as the masking
cardinality $w = |C|$ increases from $1$ (exact identification) toward $m$
(complete masking).

The purpose of this note is to show that this monotonicity is an artifact of the
uniform masking assumption, not a general feature of masked data. Specifically,
we construct a deterministic masking mechanism at the maximum masking cardinality
$w = m-1$ under which the FIM equals that of complete identification ($w=1$).
The key insight is that information loss depends on the \emph{structure} of the
masking mechanism---whether the failure cause can be recovered from the candidate
set---not on the candidate set cardinality per se.

This result has implications for diagnostic system design: a well-engineered
diagnostic that maps each failure mode to a unique signature, even one that
reports a large candidate set, incurs no information loss whatsoever for
parameter estimation.

\section{Model}
\label{sec:model}

\subsection{Exponential series systems}

Consider a series system of $m$ independent components with lifetimes
$\cv_1, \ldots, \cv_m$, where component $j$ has an exponential distribution
$\cv_j \sim \expdist(\tparam{\lambda}_j)$ with rate $\tparam{\lambda}_j > 0$.
The system lifetime is $\sv = \min(\cv_1, \ldots, \cv_m)$, and the failed
component is $\RV{K} = \arg\min_j \cv_j$.

The system lifetime is exponentially distributed with rate
$\Lambda = \sum_{j=1}^{m} \tparam{\lambda}_j$~\cite{barlow1975}. A crucial
property of the exponential case is that the failure cause and system lifetime
are independent.

\begin{proposition}[Independence of failure cause and system lifetime]
\label{prop:K_indep_S}
For exponential series systems, $\RV{K}$ and $\sv$ are independent with
$\Prob{\RV{K} = k} = \tparam{\lambda}_k / \Lambda$.
\end{proposition}

\begin{proof}
In a series system, $\Prob{\RV{K} = k \mid \sv = t} = h_k(t)/\sum_{l=1}^{m}
h_l(t)$, where $h_j$ is the hazard of component~$j$. For exponential
components $h_j(t) = \tparam{\lambda}_j$, so $\Prob{\RV{K} = k \mid \sv = t}
= \tparam{\lambda}_k / \Lambda$, which is free of $t$. Hence $\RV{K} \perp \sv$
and $\pdf_{\RV{K},\sv}(k,t) = (\tparam{\lambda}_k/\Lambda)\cdot\Lambda
\exp(-\Lambda t) = \tparam{\lambda}_k \exp(-\Lambda t)$.
\end{proof}

This independence is specific to the exponential (constant hazard) case; for
distributions with time-varying hazard rates, $\RV{K}$ and $\sv$ are generally
dependent~\cite{cox1959}.

\subsection{Masking mechanisms}

Upon observing a system failure at time $t$, a diagnostic procedure reports a
\emph{candidate set} $\cand \subseteq \{1,\ldots,m\}$ containing the failed
component $\RV{K}$. The \emph{masking mechanism} is the conditional distribution
$\Prob{\cand = c \mid \RV{K} = k}$, subject to the containment constraint
$k \in c$. We assume the mechanism depends only on $\RV{K}$ and not on $\sv$ or
on the parameter $\tparam{\vec\lambda}$.

The joint density of the observed data $(t, c)$ under a general masking
mechanism is:
\begin{equation}
\label{eq:general_joint_expanded}
\pdf(t, c \mid \tparam{\vec\lambda})
  = \sum_{k \in c} \Prob{\cand = c \mid \RV{K} = k} \cdot \tparam{\lambda}_k \exp(-\Lambda t).
\end{equation}
Here we have used the factorization
$\pdf_{\RV{K},\sv}(k,t) = \tparam{\lambda}_k \exp(-\Lambda t)$ from
Proposition~\ref{prop:K_indep_S} and the assumption that masking is independent
of $\sv$.

\subsection{Fisher information under general masking}

The FIM for the parameter $\tparam{\vec\lambda}$ from a single observation
$(t, c)$ under a general masking mechanism is:
\begin{equation}
\label{eq:general_fim}
[\infomatrx(\tparam{\vec\lambda})]_{j\ell}
  = - \Expect{\frac{\partial^2}{\partial \lambda_j \partial \lambda_\ell}
    \ln \pdf(\sv, \cand \mid \vec\lambda)}\bigg|_{\tparam{\vec\lambda}}.
\end{equation}
For \emph{uniform masking} at cardinality $w$, where
$\Prob{\cand = c \mid \RV{K} = k} = 1/\binom{m-1}{w-1}$ for each $c$ of size
$w$ containing $k$, the FIM takes the form~\cite{towell2025expo}:
\begin{equation}
\label{eq:uniform_fim}
[\infomatrx(\tparam{\vec\lambda} \mid w)]_{j\ell}
  = \frac{1}{\binom{m-1}{w-1}\,\Lambda}
    \sum_{\substack{c:\,|c|=w \\ j \in c,\; \ell \in c}}
    \frac{1}{\sum_{p \in c} \tparam{\lambda}_p}.
\end{equation}

At $w=1$ (exact identification), each observation reveals $\RV{K}$ directly, and
the FIM is diagonal:
\begin{equation}
\label{eq:fim_w1}
[\infomatrx(\tparam{\vec\lambda} \mid w\!=\!1)]_{j\ell}
  = \frac{\delta_{j\ell}}{\Lambda\, \tparam{\lambda}_j},
\end{equation}
where $\delta_{j\ell}$ is the Kronecker delta.

\section{Main Result}
\label{sec:main}

\begin{definition}[Separating-deterministic masking]
\label{def:sep_det}
Let $\binom{[m]}{w}$ denote the collection of $w$-element subsets of
$\{1,\ldots,m\}$. A masking mechanism at cardinality $w$ is
\emph{separating-deterministic} if there exists an injection
$\varphi\colon \{1,\ldots,m\} \to \binom{[m]}{w}$ satisfying:
\begin{enumerate}
\item $k \in \varphi(k)$ for each $k$ (containment), and
\item $\Prob{\cand = \varphi(k) \mid \RV{K} = k} = 1$ for each $k$
      (determinism).
\end{enumerate}
\end{definition}

The injectivity of $\varphi$ means that distinct failure causes produce distinct
candidate sets, so $\varphi^{-1}$ recovers $\RV{K}$ with certainty from the
observed $\cand$.

\begin{theorem}[Information recovery under deterministic masking]
\label{thm:main}
For an $m$-component exponential series system with rates
$\tparam{\vec\lambda} = (\tparam{\lambda}_1,\ldots,\tparam{\lambda}_m)$,
suppose the masking mechanism is separating-deterministic with injection
$\varphi$. Then the Fisher information matrix equals that of complete
identification:
\begin{equation}
\label{eq:main_result}
[\infomatrx(\tparam{\vec\lambda} \mid \varphi)]_{j\ell}
  = \frac{\delta_{j\ell}}{\Lambda\, \tparam{\lambda}_j}
  = [\infomatrx(\tparam{\vec\lambda} \mid w\!=\!1)]_{j\ell}.
\end{equation}
\end{theorem}

\begin{proof}
Under a separating-deterministic mechanism with injection $\varphi$, the masking
probabilities are $\Prob{\cand = c \mid \RV{K} = k} = \mathbbm{1}\{c =
\varphi(k)\}$. Substituting into~\eqref{eq:general_joint_expanded}:
\begin{equation}
\label{eq:det_joint}
\pdf(t, c \mid \tparam{\vec\lambda})
  = \sum_{k \in c} \mathbbm{1}\{c = \varphi(k)\} \cdot \tparam{\lambda}_k
    \exp(-\Lambda t).
\end{equation}
Since $\varphi$ is injective, for each candidate set $c$ in the image of
$\varphi$, there is exactly one $k^*(c) = \varphi^{-1}(c)$ such that
$\varphi(k^*) = c$. The sum collapses to:
\begin{equation}
\pdf(t, c \mid \tparam{\vec\lambda})
  = \tparam{\lambda}_{k^*(c)} \exp(-\Lambda t).
\end{equation}
Taking the logarithm:
\begin{equation}
\ln \pdf(t, c \mid \tparam{\vec\lambda})
  = \ln \tparam{\lambda}_{k^*(c)} - \Lambda\, t.
\end{equation}
This has the same functional form as the log-density under exact identification
($w=1$), where one observes $(t, k)$ with density
$\tparam{\lambda}_k \exp(-\Lambda t)$. Specifically, $c$ plays the role of $k$
through the bijection $k^* = \varphi^{-1}(c)$.

Computing second derivatives:
\begin{equation}
\frac{\partial^2}{\partial \lambda_j \partial \lambda_\ell}
  \ln \pdf(t, c \mid \vec\lambda)
  = - \frac{\delta_{j,k^*(c)} \cdot \delta_{\ell,k^*(c)}}{\lambda_{k^*(c)}^2}.
\end{equation}
Taking the negative expectation:
\begin{align}
[\infomatrx(\tparam{\vec\lambda} \mid \varphi)]_{j\ell}
  &= \sum_{c \in \mathrm{Im}(\varphi)} \int_0^\infty
     \frac{\delta_{j,k^*(c)} \cdot \delta_{\ell,k^*(c)}}{(\tparam{\lambda}_{k^*(c)})^2}
     \cdot \tparam{\lambda}_{k^*(c)} \exp(-\Lambda t)\, dt \nonumber\\
  &= \sum_{k=1}^{m}
     \frac{\delta_{jk}\,\delta_{\ell k}}{(\tparam{\lambda}_k)^2}
     \cdot \tparam{\lambda}_k \cdot \frac{1}{\Lambda}
  = \frac{\delta_{j\ell}}{\Lambda\, \tparam{\lambda}_j},
\end{align}
where in the second line we re-indexed the sum using $k = k^*(c) =
\varphi^{-1}(c)$ as $c$ ranges over $\mathrm{Im}(\varphi)$, and used
$\int_0^\infty e^{-\Lambda t}\, dt = 1/\Lambda$.
\end{proof}

\begin{corollary}[Existence at maximum cardinality]
\label{cor:derangement}
For $w = m-1$ and $m \geq 2$, separating-deterministic mechanisms exist.
In particular, any derangement $\sigma$ of $\{1,\ldots,m\}$ (a permutation with
no fixed points) defines such a mechanism via
$\varphi(k) = \{1,\ldots,m\} \setminus \{\sigma(k)\}$.
\end{corollary}

\begin{proof}
We verify the two conditions of Definition~\ref{def:sep_det}.
\emph{Containment:} Since $\sigma$ is a derangement, $\sigma(k) \neq k$, so
$k \notin \{\sigma(k)\}$ and hence $k \in \{1,\ldots,m\}\setminus\{\sigma(k)\}
= \varphi(k)$.
\emph{Injectivity:} If $\varphi(k) = \varphi(k')$, then
$\{1,\ldots,m\}\setminus\{\sigma(k)\} = \{1,\ldots,m\}\setminus\{\sigma(k')\}$,
so $\sigma(k) = \sigma(k')$, and since $\sigma$ is a permutation, $k = k'$.
Derangements of $\{1,\ldots,m\}$ exist for all $m \geq 2$ (the simplest is the
cyclic shift $\sigma(k) = (k \bmod m) + 1$).
\end{proof}

\begin{remark}[Diagnostic information vs.\ estimation precision]
\label{rem:IKC_vs_FIM}
Proposition~\ref{prop:K_indep_S} enables a clean separation between two
information-theoretic quantities:
\begin{itemize}
\item The \emph{mutual information} $\mathrm{I}(\RV{K}; \cand)$ measures how
  well the candidate set identifies the failed component. Under uniform masking,
  this decreases monotonically with $w$.
\item The \emph{Fisher information matrix} $\infomatrx(\tparam{\vec\lambda})$
  measures how precisely the parameters can be estimated.
\end{itemize}
These quantities need not track together. Under separating-deterministic masking
at $w = m-1$, $\cand = \varphi(\RV{K})$ is a deterministic function of
$\RV{K}$, so $\mathrm{H}(\RV{K} \mid \cand) = 0$ and
$\mathrm{I}(\RV{K}; \cand) = \mathrm{H}(\RV{K})$ (maximal
diagnostic information), and the FIM matches
$w=1$. Under uniform masking at $w = m-1$, $\mathrm{I}(\RV{K};\cand)$ is
minimized among all identifiable mechanisms, and the FIM is strictly degraded.
The separating-deterministic mechanism achieves both maximal diagnostic
information and maximal estimation precision simultaneously.

Note that the separating-deterministic mechanism violates the symmetry condition
C2 of the standard C1--C2--C3 framework~\cite{towell2023reliability}: we have
$\Prob{\cand = \varphi(k) \mid \RV{K} = k} = 1$ but
$\Prob{\cand = \varphi(k) \mid \RV{K} = k'} = 0$ for $k' \neq k$ with
$k' \in \varphi(k)$. The C2 assumption---which requires masking probabilities to
be symmetric across components within the candidate set---is precisely what costs
information under uniform masking.
\end{remark}

\section{Numerical Verification}
\label{sec:numerical}

We verify Theorem~\ref{thm:main} via Monte Carlo simulation. For each of two
configurations ($m=3$ with $\tparam{\vec\lambda} = (1,3,5)$ and $m=5$ with
$\tparam{\vec\lambda} = (1,2,3,5,7)$), we generate $R = 100{,}000$ independent
samples of size $n$ from the joint distribution $(\sv, \RV{K})$. From each
sample, we construct three datasets by applying three different masking
mechanisms to the \emph{same} underlying $(\sv, \RV{K})$ data:
\begin{enumerate}
\item \textbf{No masking} ($w=1$): observe $\RV{K}$ directly.
\item \textbf{Deterministic} ($w=m-1$): apply the cyclic derangement
  $\sigma(k) = (k \bmod m) + 1$ and set
  $\cand = \{1,\ldots,m\}\setminus\{\sigma(\RV{K})\}$.
\item \textbf{Uniform} ($w=m-1$): include each non-failed component
  independently with equal probability; equivalently, exclude one uniformly
  random non-failed component.
\end{enumerate}

For mechanisms~(1) and~(2), the MLE is
$\estimator{\lambda}_j = n_j / (n\bar{t})$ where $n_j$ counts observations
attributed to component~$j$ (for mechanism~(2), the recovery
$\varphi^{-1}(\cand)$ yields the same counts as mechanism~(1)). For
mechanism~(3), we use the closed-form MLE for uniform masking at
$w=m-1$~\cite{towell2025expo}.

Table~\ref{tab:variances} reports the empirical variances (scaled by $n$) at
$n = 2000$ alongside theoretical predictions. The quantities $n \cdot
\mathrm{Var}(\estimator{\lambda}_j)$ converge to
$[\infomatrx^{-1}]_{jj}$~\cite{bickel}. As predicted by
Theorem~\ref{thm:main}, mechanisms~(1) and~(2) yield identical empirical
variances, both matching the $w=1$ theory $\Lambda \tparam{\lambda}_j$. The
uniform mechanism produces substantially larger variances.

\begin{table}[ht]
\centering
\caption{Scaled empirical variances $n \cdot \mathrm{Var}(\estimator{\lambda}_j)$ at $n = 2000$,
  $R = 100{,}000$ replications. Theory: $\Lambda \tparam{\lambda}_j$ for $w\!=\!1$
  and deterministic.}
\label{tab:variances}
\smallskip
\small
\resizebox{\textwidth}{!}{%
\begin{tabular}{@{}llccc@{}}
\toprule
Config.\ & Mechanism & $j=1$ & $j=2$ & $j=3$ \\
\midrule
\multirow{3}{*}{$m\!=\!3$, $\vec\lambda\!=\!(1,3,5)$}
  & $w\!=\!1$ (theory / emp.) & 9.00 / 9.02 & 27.00 / 27.08 & 45.00 / 45.09 \\
  & Deterministic $w\!=\!2$   & 9.00 / 9.02 & 27.00 / 27.08 & 45.00 / 45.09 \\
  & Uniform $w\!=\!2$         & 81.0 / 80.8 & 81.0 / 81.5 & 81.0 / 80.5 \\
\midrule
\multirow{3}{*}{$m\!=\!5$, $\vec\lambda\!=\!(1,2,3,5,7)$}
  & $w\!=\!1$ (theory / emp.) & 18.00 / 18.04 & 36.00 / 36.04 & 54.00 / 54.16 \\
  & Deterministic $w\!=\!4$   & 18.00 / 18.04 & 36.00 / 36.04 & 54.00 / 54.16 \\
  & Uniform $w\!=\!4$         & 936 / 708 & 900 / 865 & 864 / 849 \\
\bottomrule
\end{tabular}}%

\medskip
\footnotesize
Note: Columns $j=4,5$ omitted for space; all components exhibit the same pattern.
For $m\!=\!5$ uniform masking, 7.9\% of replicates produced boundary MLEs
(some $\estimator{\lambda}_j \leq 0$) and were excluded, causing slower empirical
convergence; at $n = 10{,}000$ the empirical values match theory to within 1\%.
\end{table}

Figure~\ref{fig:verification} shows the convergence of
$\mathrm{Var}(\estimator{\lambda}_1)$ across sample sizes $n$ ranging from $50$
to $2000$. In both configurations, the no-masking and
deterministic curves are indistinguishable and converge to the same theoretical
prediction, while the uniform masking curve converges to a strictly larger
variance.

\begin{figure}[ht]
\centering
\includegraphics[width=\textwidth]{fig_deterministic_masking.pdf}
\caption{Empirical variance of $\estimator{\lambda}_1$ under three masking
  mechanisms. (a) $m=3$, $\tparam{\vec\lambda} = (1,3,5)$.
  (b) $m=5$, $\tparam{\vec\lambda} = (1,2,3,5,7)$. Dotted lines show
  theoretical predictions. The no-masking and deterministic curves are
  indistinguishable, confirming Theorem~\ref{thm:main}.}
\label{fig:verification}
\end{figure}

\section{Discussion}
\label{sec:discussion}

\paragraph{Practical implications.}
Theorem~\ref{thm:main} shows that a diagnostic procedure need not identify the
exact failure cause to achieve full estimation efficiency---it suffices for each
failure mode to produce a \emph{unique diagnostic signature}. In engineering
applications, a diagnostic that consistently maps failure mode $k$ to a specific
alarm pattern $\varphi(k)$, even one involving multiple components, is as
informative as exact identification for the purpose of estimating component
failure rates.

\paragraph{Limitations.}
The result depends on the exponential lifetime assumption through
Proposition~\ref{prop:K_indep_S}: the independence of $\RV{K}$ and $\sv$ is what
allows the masking mechanism to be ``undone'' by the injection $\varphi^{-1}$.
For Weibull or other distributions where $\RV{K}$ depends on $\sv$, the
conditional distribution $\pdf(t \mid \RV{K} = k)$ varies with $k$, and
observing $\cand = \varphi(k)$ recovers $\RV{K}$ but not the full conditional
structure. Whether an analogous result holds approximately when hazard rates vary
slowly is an open question.

\paragraph{Open questions.}
Two directions are natural. First, a \emph{perturbation analysis}: how does the
FIM degrade continuously as the masking mechanism moves from deterministic toward
uniform? Characterizing the functional form of the information loss as the
mechanism becomes ``noisier'' would connect the discrete result of
Theorem~\ref{thm:main} to the information-loss landscape. Second,
\emph{optimal mechanism design}: among all masking mechanisms at a given
cardinality $w$, which maximizes the determinant of the FIM? Since
$\binom{m}{w} \geq m$ for all $1 \leq w \leq m-1$, a
separating-deterministic mechanism always exists, suggesting the
optimum may always be deterministic, but this requires proof.

\section*{Acknowledgements}
This result was developed as part of ongoing work extending the author's
master's project~\cite{towell2023reliability}. Simulation code and data are
available at \url{https://github.com/queelius/deterministic-masking}.

\section*{Declaration of generative AI in the manuscript preparation process}
During the preparation of this work the author used Claude (Anthropic) in order
to assist with \LaTeX{} formatting and document class conversion. After using
this tool, the author reviewed and edited the content as needed and takes full
responsibility for the content of the published article.

\bibliographystyle{elsarticle-num}
\bibliography{references}

\end{document}
